\section{Tema, recorte e objetivos}

\subsection{Tema e t\'itulo provis\'orio}
Tema: economia comparada dos combust\'iveis f\'osseis. T\'itulo provis\'orio: \emph{Economia de Combust\'iveis F\'osseis no Brasil e no Mundo: Uma An\'alise Comparativa de Petr\'oleo, G\'as Natural e Carv\~ao Mineral}.

\subsection{Pergunta de pesquisa}
Como diferem, entre Brasil e mundo, a din\^amica de reservas, produ\c{c}\~ao, consumo, com\'ercio e pre\c{c}os de petr\'oleo, g\'as natural e carv\~ao mineral, e quais implica\c{c}\~oes econ\^omicas decorrem dessas diferen\c{c}as?

\subsection{Objetivo geral}
Comparar a economia de petr\'oleo, g\'as natural e carv\~ao mineral no Brasil e no mundo com base em indicadores padronizados e compar\'aveis.

\subsection{Objetivos espec\'ificos}
\begin{itemize}
  \item consolidar um conjunto comum de indicadores para os tr\^es combust\'iveis;
  \item analisar diferen\c{c}as de oferta, demanda e com\'ercio entre Brasil e mundo;
  \item comparar a forma\c{c}\~ao de pre\c{c}os e fatores de volatilidade;
  \item discutir implica\c{c}\~oes para seguran\c{c}a energ\'etica e planejamento setorial.
\end{itemize}

\subsection{Recorte}
\textbf{Combust\'iveis}: petr\'oleo, g\'as natural e carv\~ao mineral. \\
\textbf{Abrang\^encia geogr\'afica}: Brasil e mundo. \\
\textbf{Horizonte temporal}: s\'eries hist\'oricas recentes (prioridade para 2015--2025, conforme disponibilidade) e proje\c{c}\~oes quando explicitamente indicadas pelas fontes.

\subsection{Or\c{c}amento de conte\'udo (meta: 17--22 p\'aginas)}
\begin{table}[h]
\centering
\caption{Distribui\c{c}\~ao planejada de p\'aginas}
\begin{tabular}{lc}
\toprule
Se\c{c}\~ao & Faixa (p\'aginas) \\
\midrule
Introdu\c{c}\~ao + justificativa & 2--3 \\
Tema, recorte e m\'etodo & 2--3 \\
An\'alise do petr\'oleo & 3--4 \\
An\'alise do g\'as natural & 3--4 \\
An\'alise do carv\~ao mineral & 3--4 \\
Discuss\~ao comparativa & 2--3 \\
Conclus\~ao & 1--2 \\
\midrule
Total previsto & 17--22 \\
\bottomrule
\end{tabular}
\end{table}
